% Module 4: How to Make Qubits?
% Estimated slides: 25 | Duration: ~1.5 hours
% Key topics: Physical qubit technologies, their pros/cons,
%             fidelity, coherence times, qubit counts, company landscape
% Labs: IBM Quantum device explorer tour

%%%%%%%%%%%%%%%%%%%%%%%%%%%%%%%%%%%%%%%%%%%%%%%%%%%%%%%%%%
\begin{frame}[fragile]\frametitle{}
\begin{center}
{\Large Module 4: How to Make Qubits?}
\end{center}
\end{frame}

%%%%%%%%%%%%%%%%%%%%%%%%%%%%%%%%%%%%%%%%%%%%%%%%%%%%%%%%%%%
\begin{frame}[fragile]\frametitle{The Challenge: Quantum Objects are Fragile}
\begin{itemize}
\item Quantum superpositions collapse when they interact with the environment
\item Building a qubit means isolating a quantum object from disturbance
\item Requirements: extreme cold, electromagnetic shielding, vibration isolation
\item At the same time: we need to control and read qubits precisely
\item Competing demands: isolate from environment, but couple to control electronics
\item Different physical technologies make different trade-offs
\end{itemize}
% Suggested diagram: Qubit ``sandwich'': isolated from environment on outside, controlled from inside
\end{frame}

%%%%%%%%%%%%%%%%%%%%%%%%%%%%%%%%%%%%%%%%%%%%%%%%%%%%%%%%%%%
\begin{frame}[fragile]\frametitle{Key Qubit Quality Metrics}
\begin{itemize}
\item \textbf{T1 (energy relaxation time)}: how long a $|1\rangle$ state survives before decaying to $|0\rangle$
\item \textbf{T2 (dephasing time)}: how long phase information (superposition) survives
\item \textbf{Gate fidelity}: accuracy of single and two-qubit gates (target: >99.9\%)
\item \textbf{Readout fidelity}: accuracy of measuring a qubit state
\item \textbf{Connectivity}: which pairs of qubits can interact directly
\item Typical IBM superconducting: T1 = 50-300 microseconds, gate fidelity 99-99.9\%
\end{itemize}
% Suggested diagram: Quality metrics scorecard for different qubit technologies
\end{frame}

%%%%%%%%%%%%%%%%%%%%%%%%%%%%%%%%%%%%%%%%%%%%%%%%%%%%%%%%%%%
\begin{frame}[fragile]\frametitle{Technology 1: Superconducting Qubits}
\begin{itemize}
\item Used by: IBM, Google, Rigetti, IQM
\item Physical mechanism: superconducting circuit loops (Josephson junctions) at ~15 millikelvin
\item Pros: fabricated with chip technology (scalable), fast gates (~10-50 ns)
\item Cons: requires extreme cooling (near absolute zero); shorter coherence times
\item IBM Eagle (127 qubits), Osprey (433), Condor (1121), Heron (133)
\item Google Sycamore (53 qubits), claimed quantum supremacy in 2019
\end{itemize}
% Suggested diagram: Dilution refrigerator photo schematic; IBM chip photo
\end{frame}

%%%%%%%%%%%%%%%%%%%%%%%%%%%%%%%%%%%%%%%%%%%%%%%%%%%%%%%%%%%
\begin{frame}[fragile]\frametitle{Technology 2: Trapped Ion Qubits}
\begin{itemize}
\item Used by: IonQ, Quantinuum (Honeywell), Oxford Ionics
\item Physical mechanism: individual ions (charged atoms) suspended by laser/electromagnetic fields
\item Pros: very high fidelity gates, long coherence times (seconds)
\item Cons: slower gates (~microseconds); harder to scale to many qubits
\item IonQ Forte: 36 algorithmic qubits; Quantinuum H2: 56 qubits
\item All-to-all connectivity: any qubit can interact with any other (unlike superconducting)
\end{itemize}
% Suggested diagram: Ion trap diagram showing laser control of ions in vacuum
\end{frame}

%%%%%%%%%%%%%%%%%%%%%%%%%%%%%%%%%%%%%%%%%%%%%%%%%%%%%%%%%%%
\begin{frame}[fragile]\frametitle{Technology 3: Photonic Qubits}
\begin{itemize}
\item Used by: PsiQuantum, Xanadu, QuiX Quantum
\item Physical mechanism: photons (light particles) encode quantum information in polarization or path
\item Pros: room temperature operation, low loss, natural for networking/communication
\item Cons: photon-photon interactions are hard; two-qubit gates are probabilistic
\item PsiQuantum aims for fault-tolerant QC via photonic chips at TSMC scale
\item Xanadu PennyLane framework; Borealis: Gaussian Boson Sampling
\end{itemize}
% Suggested diagram: Photonic chip schematic with waveguides and beam splitters
\end{frame}

%%%%%%%%%%%%%%%%%%%%%%%%%%%%%%%%%%%%%%%%%%%%%%%%%%%%%%%%%%%
\begin{frame}[fragile]\frametitle{Technology 4: Neutral Atom Qubits}
\begin{itemize}
\item Used by: Pasqal, QuEra Computing, Atom Computing, Infleqtion
\item Physical mechanism: neutral atoms trapped by laser tweezers in configurable arrays
\item Pros: reconfigurable connectivity, potentially large qubit counts, long coherence
\item Cons: slower gates; readout can disturb atoms
\item QuEra Aquila: 256 qubits (analog mode); Atom Computing Phoenix: 1180 qubits
\item Great for quantum simulation and optimization problems
\end{itemize}
% Suggested diagram: Atom array photo/diagram with laser tweezers
\end{frame}

%%%%%%%%%%%%%%%%%%%%%%%%%%%%%%%%%%%%%%%%%%%%%%%%%%%%%%%%%%%
\begin{frame}[fragile]\frametitle{Technology 5: Silicon Spin Qubits}
\begin{itemize}
\item Used by: Intel, Silicon Quantum Computing (Australia)
\item Physical mechanism: electron or nuclear spin of atoms in silicon chips at cryogenic temperatures
\item Pros: leverages existing semiconductor manufacturing; small qubit size
\item Cons: still early stage; T2 times are short; difficult to read out
\item Intel Tunnel Falls: 12-qubit chip (2023); targeting millions of qubits long-term
\item Compatible with classical chip fabrication -- potential for integration
\end{itemize}
% Suggested diagram: Silicon spin qubit chip cross-section
\end{frame}

%%%%%%%%%%%%%%%%%%%%%%%%%%%%%%%%%%%%%%%%%%%%%%%%%%%%%%%%%%%
\begin{frame}[fragile]\frametitle{Technology 6: Topological Qubits}
\begin{itemize}
\item Pursued by: Microsoft (Azure Quantum)
\item Physical mechanism: Majorana fermion quasi-particles in topological materials
\item Pros: theoretically very low error rates (topological protection)
\item Cons: extremely difficult to realize; still in research phase
\item Microsoft announced first topological qubit chip (2023) -- very early stage
\item If successful: could enable fault-tolerant QC with far fewer physical qubits
\end{itemize}
% Suggested diagram: Topological qubit schematic showing Majorana modes at wire ends
\end{frame}

%%%%%%%%%%%%%%%%%%%%%%%%%%%%%%%%%%%%%%%%%%%%%%%%%%%%%%%%%%%
\begin{frame}[fragile]\frametitle{Technology Comparison}
\begin{itemize}
\item No single technology wins on all metrics --- different use cases favor different approaches
\end{itemize}
\begin{center}
% Suggested diagram: Comparison table with rows: Technology | Company | Qubit count | T2 | Gate speed | Fidelity | Scalability
% Superconducting | IBM/Google | 100-1000+ | microsec | fast | 99.9% | high potential
% Trapped ion | IonQ/Quantinuum | 10-100 | seconds | slow | 99.99% | medium
% Photonic | PsiQuantum | many | N/A | fast | --- | high potential
% Neutral atom | QuEra/Pasqal | 100-1000 | seconds | medium | 99%+ | high potential
% Silicon spin | Intel | 10s | short | fast | medium | high potential
\end{center}
\end{frame}

%%%%%%%%%%%%%%%%%%%%%%%%%%%%%%%%%%%%%%%%%%%%%%%%%%%%%%%%%%%
\begin{frame}[fragile]\frametitle{The Path to Fault-Tolerant QC}
\begin{itemize}
\item Current NISQ devices: hundreds to thousands of noisy physical qubits
\item Fault-tolerant QC requires logical qubits (error-corrected)
\item Each logical qubit may require 1000+ physical qubits for error correction
\item For Shor's algorithm to break RSA: ~4000 logical qubits = millions of physical qubits
\item Timeline estimates vary: 2030-2040 for practical fault-tolerant machines
\item NISQ devices are already useful today for simulation, optimization
\end{itemize}
% Suggested diagram: Road from physical qubit -> logical qubit -> fault-tolerant QC
\end{frame}

%%%%%%%%%%%%%%%%%%%%%%%%%%%%%%%%%%%%%%%%%%%%%%%%%%%%%%%%%%%
\begin{frame}[fragile]\frametitle{The Quantum Refrigerator: Inside a Dilution Refrigerator}
\begin{itemize}
\item Superconducting qubits require cooling to 15 millikelvin (~$-273.1^\circ$C)
\item Colder than outer space (which is 2.7 Kelvin = $-270.4^\circ$C)
\item Achieved with a dilution refrigerator: layers of cooling from room temp to base temp
\item Looks like a chandelier of stacked metal plates inside a cryostat
\item IBM's dilution refrigerators: $\sim$3 meters tall
\item The qubit chip is at the very bottom, isolated from all vibration and EM noise
\end{itemize}
% Suggested diagram: Dilution refrigerator cross-section showing temperature stages
\end{frame}

%%%%%%%%%%%%%%%%%%%%%%%%%%%%%%%%%%%%%%%%%%%%%%%%%%%%%%%%%%%
\begin{frame}[fragile]\frametitle{Lab: IBM Quantum Device Explorer}
\begin{itemize}
\item Access IBM Quantum at \texttt{quantum.ibm.com}
\item Explore the list of available quantum systems
\item View device topology: qubit connectivity graph
\item View calibration data: T1, T2, gate fidelities per qubit
\item Observe how calibration data changes day to day
\item Identify which qubits have the best T2 and gate fidelity
\end{itemize}
\begin{verbatim}
from qiskit_ibm_runtime import QiskitRuntimeService
service = QiskitRuntimeService()
backend = service.backend("ibm_brisbane")
print(backend.properties())
\end{verbatim}
\end{frame}

%%%%%%%%%%%%%%%%%%%%%%%%%%%%%%%%%%%%%%%%%%%%%%%%%%%%%%%%%%%
\begin{frame}[fragile]\frametitle{Module 4 Summary}
\begin{itemize}
\item Multiple physical qubit technologies exist; all are qubits abstractly
\item Key metrics: T1/T2 coherence times, gate fidelity, connectivity, qubit count
\item Superconducting (IBM, Google): fastest gates, most scalable currently
\item Trapped ion (IonQ, Quantinuum): highest fidelity, all-to-all connectivity
\item Neutral atom, photonic, silicon spin: emerging alternatives
\item NISQ era: noisy devices useful now; fault-tolerant QC is the goal
\item Dilution refrigerators required for superconducting qubits (15 mK)
\end{itemize}
\end{frame}
